\scanpage{291}
\uline{Structure de certains groupes divisibles}.

\deleted{\unsure{Les}} Les groupes de torsion divisibles $\unsure{\mathfrak M|M}$ sont comme l'on sait isomorphes à des sommes directes de groupes $\mathbb Q_p/\mathbb Z_p$ ($p$ premiers); Ceux qui sont $p$-primaires sont donc isom. à des \deleted{sommes directes de} groupes $(\mathbb Q_p/\mathbb Z_p)^{(I)}$, où le cardinal de $I$ est bien déterminé; c'est aussi la dimension de l'espace vectoriel sur $\unsure{p}$ ${}_p\mathfrak M=\operatorname{Hom}(\mathbb Z/p\mathbb Z,\mathfrak M)$.
\deleted{\unsure{Les $\mathfrak M$ en question} de torsion $p$-primaires sont d'ailleurs isomorphes}

Les conditions équivalentes

(i) $\mathfrak M$ groupe de torsion $p$-primaire, et ${}_p\mathfrak M$ est fini

(ii) $\mathfrak M$ isomorphe à un sous-groupe d'un $(\mathbb Q_p/\mathbb Z_p)^{\unsure{n}}$, \deleted{\unsure{$n\in\mathbb Z$}} $n$ entier $\geq0$.

(iii) $\mathfrak M$ isomorphe à un groupe $F+(\mathbb Q_p/\mathbb Z_p)^{\rho}$, où $F$ est un $p$-groupe de torsion fini, et $\rho$ un entier fini.
\margin{Modules sur un anneau de Dedekind.}
Sous ces conditions, $\rho$ est défini de façon unique [c'est aussi $\operatorname{rang}_{\mathbb Z/p\mathbb Z}{}_p\mathfrak M^0$, où $\mathfrak M^0$ est le plus grand sous-groupe divisible de $\mathfrak M$] ainsi que la structure de $F$ [c'est $\mathfrak M/\mathfrak M^0$].

On peut de deux façons différentes faire correspondre à $\mathfrak M$ un module de type fini sur $\mathbb Z_p$:
On peut poser
\[
T_p(\mathfrak M)=\operatorname{Hom}(\mathfrak M,\mathbb Q_p/\mathbb Z_p)=\varprojlim_n\operatorname{Hom}({}_{p^n}\mathfrak M,\mathbb Q_p/\mathbb Z_p)
\]
\margin{$T_p$ \uline{foncteur} \uline{contravariant} en $M$}
[$\mathfrak M\mapsto T_p(\mathfrak M)$ établit une \inserted{anti-}équivalence de la catégorie des $\mathfrak M$ satisfaisant (i), et des modules de type fini sur $\mathbb Z_p$.

On peut prendre aussi
\[
T^p(M)=\varprojlim {}_{p^n}M
\]
\margin{$T^p$ foncteur \uline{covariant} en $M$}
où pour $m\geq n$, on envoie ${}_{p^m}M$ dans ${}_{p^n}M$ par $p^{m-n}$.
\figure{\begin{tikzpicture}\draw (0,0)..controls (0.7,-0.08) and (1.5,0.05)..(2.1,0.09)..controls (2.4,0.1) and (2.25,0.65)..(2.4,0.68)--(4.1,0.67);\end{tikzpicture}}
\mspage{290}
